\documentclass[11pt]{ctexart}
\usepackage[a4paper,margin=1in]{geometry}
\usepackage{graphicx}
\usepackage{xcolor}
\usepackage{hyperref}
\definecolor{refgray}{gray}{0.45}
\title{\textbf{市场最新观点汇总}\\[0.4em]{\large 全球FDI流向战略与安全，政策不确定性拖累清洁转型，零售业供应链与定价面临AI重塑}}
\author{KC桌面 自动生成}
\date{260708}
\begin{document}
\maketitle
\tableofcontents
\newpage
\section{一页摘要}
\begin{itemize}
  \item 全球FDI在2025年增长6\%至1.6万亿美元，但增长高度集中于AI、半导体等战略行业，发展中国家被边缘化。
  \item 亚洲开发银行研究显示，政策不确定性每上升一个标准差，新兴市场碳排放增加0.03\%，碳定价和金融发展是有效缓冲。
  \item 波士顿咨询和麦肯锡的多份报告共同指向零售业和技能市场的结构性变革：AI定价、SKU精简、自有品牌战略和供应链协同成为竞争关键。
\end{itemize}
\section{投行/券商}
今日暂无新增报告。

\begin{itemize}
  \item 今日暂无新增报告。
\end{itemize}
\section{战略咨询}
波士顿咨询和麦肯锡聚焦零售业供应链、定价、SKU管理及AI对就业市场的影响，揭示行业结构性变革。

\subsection{零售供应链与运营效率：从博弈到协同}
\textbf{生鲜供应链的瓶颈在于零售商与供应商之间的信息孤岛和激励错位，而非技术本身。}

\begin{itemize}
  \item 波士顿咨询指出，生鲜SKU增加导致出库卡车装载率下降、紧急补货增多、损耗率上升，零售商与CPG供应商的博弈关系已失效，必须转向深度数字化协同。
  \item SKU膨胀是传统超市的隐形杀手：过去30年SKU数量增长80\%，但试验表明削减20\%-25\%的SKU后，受影响品类销售额反而增长2\%，因为消费者面临“选择瘫痪”。
  \item 年轻家庭和高收入夫妇的线上杂货消费比纯线下高出30\%-50\%，但零售商供给端犹豫是市场爆发的主要制约，建议优先发展“点击取货”模式。
\end{itemize}
\begin{figure}[htbp]
\centering
\includegraphics[width=0.88\linewidth]{figures/fig_181.png}
\caption{波士顿咨询视觉摘要 1 - 波士顿咨询｜波士顿咨询：生鲜供应链的瓶颈不是技术，而是组织协同｜用于快速识别该外部报告的核心问题和市场含义。}
\end{figure}
\begin{figure}[htbp]
\centering
\includegraphics[width=0.88\linewidth]{figures/fig_182.png}
\caption{波士顿咨询视觉摘要 1 - 波士顿咨询｜波士顿咨询：超市SKU不是太少，而是太多，减20\%反增销售额｜用于快速识别该外部报告的核心问题和市场含义。}
\end{figure}
\begin{figure}[htbp]
\centering
\includegraphics[width=0.88\linewidth]{figures/fig_190.png}
\caption{波士顿咨询视觉摘要 1 - 波士顿咨询｜波士顿咨询：年轻家庭线上杂货消费比线下高30\%至50\%｜用于快速识别该外部报告的核心问题和市场含义。}
\end{figure}
{\small\color{refgray}\textbf{References:} [R010] 波士顿咨询：生鲜供应链的瓶颈不是技术，而是组织协同; [R011] 波士顿咨询：超市SKU不是太少，而是太多，减20\%反增销售额; [R012] 波士顿咨询：年轻家庭线上杂货消费比线下高30\%至50\%}

\subsection{AI与定价策略：从规则驱动到智能动态定价}
\textbf{AI驱动的动态定价已成为零售商维持竞争力的必要条件，可同时处理战略、卫生和动态三大维度。}

\begin{itemize}
  \item 已落地AI定价的零售商毛利润提升5\%-10\%，营收和客户价值感知同步增长，核心在于“去平均化”定价而非简单调价。
  \item 传统定价只关注成本加成或竞品对标，AI引擎可同时考虑十几个变量，在数十亿场景中找到每个门店、每个商品的最优价格。
  \item 动态定价的核心不是降价，而是精准定位：一家连锁超市通过系统追踪竞品，发现部分商品价格低于竞品20\%-30\%，提价后利润改善。
\end{itemize}
\begin{figure}[htbp]
\centering
\includegraphics[width=0.88\linewidth]{figures/fig_191.jpg}
\caption{Exhibit 1 - Exhibit 1 - The Impact of Improving Pricing Capabilities With AI-powered solutions, retailers can translate their strategic choices into optimal prices for each product and store. They can respond dynamically to both internal and}
\end{figure}
\begin{figure}[htbp]
\centering
\includegraphics[width=0.88\linewidth]{figures/fig_192.jpg}
\caption{Exhibit 2 - Exhibit 2 - AI-Powered Engines Embrace the Full Complexity of Retail Pricing The dimensions of an AI-powered approach to pricing fall into three categories: strategic, hygienic, and dynamic. AI-powered solutions can iterate throug}
\end{figure}
\begin{figure}[htbp]
\centering
\includegraphics[width=0.88\linewidth]{figures/fig_193.png}
\caption{波士顿咨询视觉摘要 1 - 波士顿咨询｜波士顿咨询：AI定价不是选择题，是零售商的生存题｜用于快速识别该外部报告的核心问题和市场含义。}
\end{figure}
{\small\color{refgray}\textbf{References:} [R013] 波士顿咨询：AI定价不是选择题，是零售商的生存题}

\subsection{自有品牌战略与技能市场重塑}
\textbf{自有品牌已从低价仿制品进化为零售商差异化、提升利润的核心武器；AI匹配工具正改写技能培训与就业规则。}

\begin{itemize}
  \item 美国自有品牌年销售额达1200亿美元，占市场18\%，份额转移几乎完全由自有品牌驱动，传统超市若不加速布局将失去市场地位。
  \item 千禧一代更看重“价值感”而非绝对低价，自有品牌通过独家产品和品质对标全国品牌，能同时满足价值与差异化需求。
  \item 麦肯锡报告显示，AI教练工具已帮助50万以上学生提升表达能力，技能重塑的规模效应正在改写劳动力市场规则。
\end{itemize}
\begin{figure}[htbp]
\centering
\includegraphics[width=0.88\linewidth]{figures/fig_194.jpg}
\caption{EXHIBIT 1 - As we look at the overall rise of discount and other formats in the US, the correlation is clear: traditional supermarkets are in decline, and channels that rely more heavily on private brands are gaining share. Private brands used to be nothing more than lowe}
\end{figure}
\begin{figure}[htbp]
\centering
\includegraphics[width=0.88\linewidth]{figures/fig_195.jpg}
\caption{EXHIBIT 1 - CONSUMERS ARE SHIFTING FROM SUPERMARKETS TO WHOLESALE CLUB, CONVENIENCE, OR DISCOUNT STORES Sources: Planet Retail; Nielsen/PLMA 2018 Yearbook; BCG analysis. THE SHARE OF PRIVATE-LABEL PRODUCTS IS GROWING AT MASS, CLUB, AND DISCOUNT STORES AS CONSUMERS SEEK DI}
\end{figure}
\begin{figure}[htbp]
\centering
\includegraphics[width=0.88\linewidth]{figures/fig_196.jpg}
\caption{Exhibit 2 - \textbackslash{}- Texas retailer HEB has become a destination for innovative private-label products, which now make up 28\% of its online sales. The company even promoted its brand during the 2019 Super Bowl. \textbackslash{}- In Europe, Tesco and Carrefour have formed a buying alliance to }
\end{figure}
{\small\color{refgray}\textbf{References:} [R014] 波士顿咨询：自有品牌不是平替，而是零售商核心竞争力; [R015] 麦肯锡：AI匹配工具正改写技能培训与就业市场规则}

\section{智库/国际机构}
IMF、亚开行、联合国贸发会议、WTO等机构聚焦政策不确定性、气候韧性、FDI结构性变化及贸易环境新规则。

\subsection{政策不确定性与气候转型：新兴市场的排放韧性}
\textbf{政策不确定性显著推高新兴市场碳排放，但碳定价、金融发展和政治稳定能有效缓冲冲击。}

\begin{itemize}
  \item 亚洲开发银行研究显示，在13个新兴市场33年数据中，政策不确定性每上升一个标准差，碳排放增加0.03\%，企业推迟清洁技术投资是主因。
  \item 有碳定价机制、金融体系更发达、贸易开放度更高、政治更稳定的地区，对政策不确定性的“排放敏感度”显著更低。
  \item 高收入新兴市场在面对政策不确定性时碳排放几乎无响应，而低收入国家排放显著上升且持续时间更长。
\end{itemize}
\begin{figure}[htbp]
\centering
\includegraphics[width=0.88\linewidth]{figures/fig_001.jpg}
\caption{Figure 2 - Figure 2: Dynamic Response of Carbon Emissions to Economic Policy Uncertainty, Adding the Quadratic Term of Gross Domestic Product Growth Notes: The figure plots the impulse responses of carbon emissions to a 1-standard deviatio}
\end{figure}
\begin{figure}[htbp]
\centering
\includegraphics[width=0.88\linewidth]{figures/fig_024.png}
\caption{亚洲开发银行视觉摘要 1 - 亚洲开发银行｜亚洲开发银行：政策不确定性延缓清洁转型，新兴市场如何提升排放韧性｜用于快速识别该外部报告的核心问题和市场含义。}
\end{figure}
{\small\color{refgray}\textbf{References:} [R001] 亚洲开发银行：政策不确定性延缓清洁转型，新兴市场如何提升排放韧性}

\subsection{全球FDI结构性转变：战略行业集中与边缘化风险}
\textbf{全球FDI增长由AI、半导体等战略行业驱动，但发展中国家被边缘化，投资逻辑从成本效率转向战略安全。}

\begin{itemize}
  \item 2025年全球FDI增长6\%至1.6万亿美元，但前20大东道国吸纳超80\%流入，半导体、AI、清洁能源、关键矿产四类战略行业占绿地项目近一半。
  \item 最不发达国家仅吸引不到10\%的战略行业项目，而在其他行业这一比例超20\%，新投资逻辑“奖励资金雄厚和已建立生态系统的地区”。
  \item 供应链重构带来替代制造基地和关键矿物加工中心的机遇，但基础设施、技术和人才门槛将多数低收入国家排除在外。
\end{itemize}
\begin{figure}[htbp]
\centering
\includegraphics[width=0.88\linewidth]{figures/fig_179.png}
\caption{联合国贸发会议视觉摘要 1 - 联合国贸发会议｜联合国贸发会议：AI与半导体驱动全球FDI，但发展中国家被落下｜用于快速识别该外部报告的核心问题和市场含义。}
\end{figure}
{\small\color{refgray}\textbf{References:} [R008] 联合国贸发会议：AI与半导体驱动全球FDI，但发展中国家被落下}

\subsection{小岛国经济韧性：财政紧缩、气候风险与治理挑战}
\textbf{IMF对塞拉利昂和所罗门群岛的评估显示，财政紧缩成效显著，但外部冲击和治理脆弱性使可持续性面临考验。}

\begin{itemize}
  \item 塞拉利昂国内基础收支改善7个GDP百分点，通胀从13.8\%降至4.3\%，但中东战争推高燃料成本，2026年增长预计放缓至4.0\%，通胀反弹至11.6\%。
  \item 所罗门群岛2025年GDP增长3.5\%，经常账户盈余5.6\%，但2026年预计转为赤字3.4\%，财政空间被侵蚀，政府现金储备已见底。
  \item 所罗门群岛面临海平面上升和渔业资源下降的长期风险：可开发鱼类生物量到2050年将下降约11\%，海平面上升年均损失可达GDP的0.1\%-0.4\%。
\end{itemize}
\begin{figure}[htbp]
\centering
\includegraphics[width=0.88\linewidth]{figures/fig_073.png}
\caption{IMF视觉摘要 1 - IMF｜IMF：塞拉利昂的紧缩见效了，但真正的考验在2028之前｜用于快速识别该外部报告的核心问题和市场含义。}
\end{figure}
\begin{figure}[htbp]
\centering
\includegraphics[width=0.88\linewidth]{figures/fig_111.png}
\caption{IMF视觉摘要 1 - IMF｜IMF：不是周期而是治理，IMF评估所罗门群岛政策执行能力｜用于快速识别该外部报告的核心问题和市场含义。}
\end{figure}
\begin{figure}[htbp]
\centering
\includegraphics[width=0.88\linewidth]{figures/fig_129.png}
\caption{IMF视觉摘要 1 - IMF｜IMF：不是台风而是海洋，IMF评估所罗门群岛两大长期不确定因素｜用于快速识别该外部报告的核心问题和市场含义。}
\end{figure}
{\small\color{refgray}\textbf{References:} [R003] IMF：塞拉利昂的紧缩见效了，但真正的考验在2028之前; [R004] IMF：不是周期而是治理，IMF评估所罗门群岛政策执行能力; [R005] IMF：不是台风而是海洋，IMF评估所罗门群岛两大长期不确定因素}

\subsection{贸易与环境新规则：从原则到操作}
\textbf{WTO和IMF报告显示，贸易与环境政策正从共识走向可执行框架，碳计量、绿色补贴和跨境支付成为新焦点。}

\begin{itemize}
  \item WTO TESSD五年报告指出，79个成员已从“要不要做”进入“怎么做”阶段，碳定价、边境碳调整、标准认证等六大类政策设计存在显著差异。
  \item IMF研究显示，海湾合作委员会与高加索和中亚地区贸易潜力巨大，物流和治理改善可撬动高达150\%的贸易增长，但政策环境是主要瓶颈。
  \item 基里巴斯案例揭示社会保障扩张的“双刃剑”：贫困率从21.9\%降至5.5\%，但新增转移支付中相当比例流向酒精、烟草等不利于人力资本积累的消费。
  \item 另有 1 篇同来源报告纳入 references，用于校准该来源板块覆盖面。
\end{itemize}
\begin{figure}[htbp]
\centering
\includegraphics[width=0.88\linewidth]{figures/fig_175.png}
\caption{IMF视觉摘要 1 - IMF｜IMF：不是缺贸易意愿，是政策卡住了，IMF看海湾与中亚合作｜用于快速识别该外部报告的核心问题和市场含义。}
\end{figure}
\begin{figure}[htbp]
\centering
\includegraphics[width=0.88\linewidth]{figures/fig_180.png}
\caption{世界贸易组织视觉摘要 1 - 世界贸易组织｜世界贸易组织：光伏电池到绿氢电解槽，环境产品贸易壁垒被逐一梳理｜用于快速识别该外部报告的核心问题和市场含义。}
\end{figure}
\begin{figure}[htbp]
\centering
\includegraphics[width=0.88\linewidth]{figures/fig_025.png}
\caption{IMF视觉摘要 1 - IMF｜IMF：基里巴斯案例，当社保支出增加，教育投入反而减少？｜用于快速识别该外部报告的核心问题和市场含义。}
\end{figure}
{\small\color{refgray}\textbf{References:} [R006] IMF：不是缺贸易意愿，是政策卡住了，IMF看海湾与中亚合作; [R009] 世界贸易组织：光伏电池到绿氢电解槽，环境产品贸易壁垒被逐一梳理; [R002] IMF：基里巴斯案例，当社保支出增加，教育投入反而减少？; [R007] 经合组织：经合组织支招阿根廷，NCP重建信任需先修复基础}

\section{结语}
今日国际信源主线清晰：全球资本与贸易正被战略安全和环境规则重塑，新兴市场和小岛国面临结构性挑战，而零售业和技能市场则迎来AI驱动的效率革命。
\newpage
\section{报告覆盖清单}
\begin{itemize}
  \item [R001] 智库/国际机构 | 智库/国际机构 / 政策不确定性与气候转型：新兴市场的排放韧性 - 亚洲开发银行：政策不确定性延缓清洁转型，新兴市场如何提升排放韧性
  \item [R002] 智库/国际机构 | 智库/国际机构 / 贸易与环境新规则：从原则到操作 - IMF：基里巴斯案例，当社保支出增加，教育投入反而减少？
  \item [R003] 智库/国际机构 | 智库/国际机构 / 小岛国经济韧性：财政紧缩、气候风险与治理挑战 - IMF：塞拉利昂的紧缩见效了，但真正的考验在2028之前
  \item [R004] 智库/国际机构 | 智库/国际机构 / 小岛国经济韧性：财政紧缩、气候风险与治理挑战 - IMF：不是周期而是治理，IMF评估所罗门群岛政策执行能力
  \item [R005] 智库/国际机构 | 智库/国际机构 / 小岛国经济韧性：财政紧缩、气候风险与治理挑战 - IMF：不是台风而是海洋，IMF评估所罗门群岛两大长期不确定因素
  \item [R006] 智库/国际机构 | 智库/国际机构 / 贸易与环境新规则：从原则到操作 - IMF：不是缺贸易意愿，是政策卡住了，IMF看海湾与中亚合作
  \item [R007] 智库/国际机构 | 智库/国际机构 / 贸易与环境新规则：从原则到操作 - 经合组织：经合组织支招阿根廷，NCP重建信任需先修复基础
  \item [R008] 智库/国际机构 | 智库/国际机构 / 全球FDI结构性转变：战略行业集中与边缘化风险 - 联合国贸发会议：AI与半导体驱动全球FDI，但发展中国家被落下
  \item [R009] 智库/国际机构 | 智库/国际机构 / 贸易与环境新规则：从原则到操作 - 世界贸易组织：光伏电池到绿氢电解槽，环境产品贸易壁垒被逐一梳理
  \item [R010] 战略咨询 | 战略咨询 / 零售供应链与运营效率：从博弈到协同 - 波士顿咨询：生鲜供应链的瓶颈不是技术，而是组织协同
  \item [R011] 战略咨询 | 战略咨询 / 零售供应链与运营效率：从博弈到协同 - 波士顿咨询：超市SKU不是太少，而是太多，减20\%反增销售额
  \item [R012] 战略咨询 | 战略咨询 / 零售供应链与运营效率：从博弈到协同 - 波士顿咨询：年轻家庭线上杂货消费比线下高30\%至50\%
  \item [R013] 战略咨询 | 战略咨询 / AI与定价策略：从规则驱动到智能动态定价 - 波士顿咨询：AI定价不是选择题，是零售商的生存题
  \item [R014] 战略咨询 | 战略咨询 / 自有品牌战略与技能市场重塑 - 波士顿咨询：自有品牌不是平替，而是零售商核心竞争力
  \item [R015] 战略咨询 | 战略咨询 / 自有品牌战略与技能市场重塑 - 麦肯锡：AI匹配工具正改写技能培训与就业市场规则
\end{itemize}
\section{Disclaimer}
{\scriptsize\color{refgray}
This community is a paid learning and information-sharing community created for general educational purposes only. The membership fee is charged solely for access to community discussions, educational content, organization, and operational costs. It is not an advisory fee, management fee, performance fee, brokerage fee, or compensation for personalized financial, investment, legal, tax, or professional advice.\par
I participate and share information solely as an individual and for educational discussion among community members. I am not acting as a financial adviser, investment adviser, broker, dealer, portfolio manager, legal adviser, tax adviser, accountant, fiduciary, or any other licensed professional adviser. Nothing shared in this community constitutes, and should not be interpreted as, financial advice, investment advice, legal advice, tax advice, a recommendation, solicitation, offer, endorsement, instruction, or guarantee to buy, sell, hold, short, trade, or invest in any security, cryptocurrency, fund, derivative, strategy, or other financial product.\par
All content shared in this community, including messages, discussions, links, files, charts, examples, market commentary, watchlists, AI-generated or AI-polished summaries, and third-party materials, is general, impersonal, and educational in nature. It does not take into account any individual's financial situation, investment objectives, risk tolerance, investment horizon, liquidity needs, tax status, legal restrictions, or suitability requirements. No content should be relied upon as suitable for any specific person, account, portfolio, or financial decision.\par
Information shared in this community may be incomplete, inaccurate, outdated, speculative, AI-generated, AI-edited, or based on third-party internet sources that have not been independently verified. I make no representation or warranty as to the accuracy, completeness, timeliness, reliability, or suitability of any information shared.\par
Financial markets involve substantial risk, including the possible loss of principal. Past performance, historical data, hypothetical examples, simulated results, backtests, personal opinions, or educational examples do not guarantee future results. Each member is solely responsible for conducting their own research, exercising independent judgment, and making their own decisions. Before making any financial, investment, legal, or tax decision, members should consult qualified and licensed professionals.\par
Participation in this community, payment of any membership fee, or communication with me or other members does not create any adviser-client, fiduciary, professional, contractual, agency, partnership, employment, or other advisory relationship. By joining or participating in this community, you acknowledge and agree that you are solely responsible for your own decisions, actions, transactions, profits, losses, risks, and consequences, and that I shall not be liable for any direct, indirect, incidental, consequential, financial, legal, tax, or other loss or damage arising from or related to any content, discussion, information, or material shared in this community.\par
}
\end{document}
